\section{Appendix D. Data Lineage, Access, and Reproduction}
\phantomsection
\label{app:d}
\label{sec:extension-provenance}
\label{sec:display-archive-r14}

Data Availability and Table~\ref{tab:component-access} summarize the released materials and access conditions. The \href{https://github.com/AlfredJamesLi/chinese-skillspan-benchmark/tree/main/reproduction/major_revision_20261004}{supplementary audit package} provides supporting calculations and archived experiment records.

\begin{samepage}
The terms \texttt{Gold150}, \texttt{QA150}, and \texttt{Silver-plus} denote the human reference, additional review sample, and original silver candidate pool, respectively. The main silver-supervised baselines use model-specific training/development subsets of this pool with revised labels (\hyperref[sec:jobbert-label-versions]{Appendix C}).\par
\end{samepage}

\subsection{Relationship to the Earlier Preprint}
The present article substantially revises the resource definitions, annotation procedures, and evaluations in the earlier preprint~\citep{li2026chineseskillspanpreprint}. The sample list and original annotations for the earlier 100-sample agreement report have not been recovered, so those results are not used to establish reliability for the present reference. The available experiment records do not substantiate the earlier descriptions of linking spans to concepts or evaluating on industries and time periods excluded from training. The current article therefore makes neither claim. Model results use the reference labels, supervision sets, and scoring procedures documented for each experiment.

\subsection{Source Categories and Data Preparation}
The source corpus was originally partitioned into 17,460 training, 2,143 development, and 3,237 test sentences, with 1,600, 200, and 200 source-document identifiers, respectively. These document-disjoint partitions describe the source corpus; the original test partition is not a human-annotated evaluation set in this study. Tianchi material came from the Recruitment Dataset (ID 163746). The public-institution crawler revisited selected pages weekly or monthly.

The initial annotation collection contains 2,601 records: 150 human-reference sentences and 2,451 original silver records. The silver pool comprises 100 initially reviewed records, 730 additional records from the model-disagreement queue, and 1,621 historical extraction records. The 980-record queue supplied the 150-sentence human reference and the first 830 silver records. Model-specific inclusion rules produce the supervision splits described in \nameref{sec:experiments}.

The corpus combines four batches collected through MacroData, public-institution websites, and Tianchi. The original training partition contains 10,312 graduate-recruitment sentences and 7,148 sentences from the AI collection; development contains 2,143 sentences from the AI collection; test contains 1,423 from the AI collection, 473 from Tianchi/cloud, and 1,341 from public institutions. The expanded data use three source categories: Tianchi/cloud, public institutions, and listed companies. The blinded sample contains 26, 16, and 8 sentences from these respective categories.

MacroData's operator is Chongqing Maherui Information Technology Co., Ltd. Public-institution records retain source names, titles, URLs, listed publication dates, and geographic fields where available. Publication dates are not collection dates; a complete crawl-window and source-exclusion history has not been recovered. The authors confirm academic-research use permission. Repository notices distinguish author-created annotation rights from third-party text restrictions.

Table~\ref{tab:source-stages} gives expanded-pool candidate and retained counts. Intermediate preparation counts are documented in the \href{https://github.com/AlfredJamesLi/chinese-skillspan-benchmark/tree/main/reproduction/reviewer_revision_20261005}{preparation memo}.

\begin{table}[!htbp]
\centering\small
\caption{Recovered sampling-stage counts and missing census information. These are sentence counts, not advertisement counts or mutually exclusive industry totals.}
\label{tab:source-stages}
\begin{tabularx}{\linewidth}{@{}P{.25\linewidth}P{.26\linewidth}L@{}}
\toprule Stage & Available counts & Sampling frame / unavailable detail \\
\midrule
Original four collections & 22,840; split 17,460/2,143/3,237 & AI, Cloud, public-sector, graduate recruitment. Per-source ad counts, time windows, attrition and complete document mapping not recovered. \\
Targeted study lineage & 980 conflict-origin + 1,621 other records & 150 reference + 2,451 Silver-plus. This is a selected lineage, not a probability sample. \\
New expansion candidates & Cloud 3,174; public-sector 2,381; listed-company 2,381 & Requirement-keyword-positive sentences; fixed hash ranking. Earlier B2 contributes 2,064 retained unique sentences. \\
Replacement channel subset & 2,500 candidates; alternative retains 2,459; Codex retains 2,353 & Respective exclusions: 41 (1.6\%) and 147 (5.9\%). A reason-coded attrition breakdown is not established. \\
Additional QA150 & 60 / 45 / 45 source quotas & Sampled from new Silver; machine suggestions visible. Does not extend the independent test population. \\
\bottomrule
\end{tabularx}
\end{table}


New sentences were selected using requirement-related keywords, a fixed hash-based order for selecting candidates, and quotas for each source. Before annotation, the texts were cleaned for HTML markup and whitespace, split into sentences, and filtered by length. Checks for excluding previously used texts relied on the stored sentences and available pretraining data. No checksum was retained for the 3.2-million-line pretraining file, so the data seen during pretraining cannot be fully reconstructed. Earlier annotations have been preserved; the available records identify the prompts and handbook versions for each documented batch (\hyperref[sec:guide-version-history]{Appendix A.4}).

\subsection{Split Membership and Document Overlap}
The later experiments divide the data by sentence, whereas the original 22,840-sentence corpus was divided so that each source document appeared in only one split. Original records identify source documents; records added during expansion identify source files and rows. The two identifier systems have not been linked, so shared sources can be checked only within each system. Qwen and the encoders use 2,150 original silver records for training and 169 for development.

\begin{table}[htbp]
\centering\small
\caption{Shared source-document groups in the later experimental splits.}
\label{tab:document-overlap}
\begin{tabularx}{\linewidth}{@{}Lrrr@{}}
\toprule Split comparison & Shared groups & \makecell{Records in\\first split} & \makecell{Records in\\second split} \\
\midrule
Silver train / development & 1 & 4 & 27 \\
Silver train / human reference & 59 & 531 & 143 \\
Silver development / human reference & 0 & 0 & 0 \\
Adopted expanded train / validation & 185 & 1,060 & 208 \\
Adopted expanded train / test & 188 & 1,076 & 210 \\
Alternative expanded train / validation & 186 & 1,069 & 210 \\
Alternative expanded train / test & 189 & 1,079 & 211 \\
\bottomrule
\end{tabularx}
\par\smallskip\footnotesize\RaggedRight
Counts refer to sentence records in shared source-document groups; records with different identifiers are counted separately even when their text is identical. Split order follows the first column.
\end{table}

We checked for repeated sentence texts across splits by comparing complete texts as stored and after Unicode NFC normalization, which gives equivalent Unicode character sequences a consistent representation. Neither comparison found matching texts between the original silver training and development sets used by Qwen and the encoders, or between training and the human reference.

A broader check used Unicode NFKC normalization, which also standardizes forms such as full-width and half-width characters, then removed whitespace and ignored letter case. This check found no matching complete texts between the original silver training set and the human reference. It found two training/development matches in the original silver data and one training/validation match in each expanded pool, with no training/test matches in either expanded pool. These checks compare texts without changing the stored sentences or annotation offsets; details are provided in the \href{https://github.com/AlfredJamesLi/chinese-skillspan-benchmark/blob/24847748b314b205e9b896c1f30ecc5d7c7bdf1e/reproduction/major_revision_20261004/data_audit/audit.json}{text-overlap analysis}. Table~\ref{tab:document-overlap} reports shared source-document groups separately.

The silver training set used by Qwen and the encoders contains 2,150 records representing 1,910 distinct complete texts. Records with identical text were retained as separate training examples. To assess the effects of this repetition and of shared source documents, further experiments would need to compare training with and without repeated texts and evaluate models on sentences from documents excluded from training.

\paragraph{Blinded-sample overlap.}
The \href{https://github.com/AlfredJamesLi/chinese-skillspan-benchmark/tree/main/reproduction/reviewer_supplement_20261004}{sample-overlap check} found no shared sample IDs or matching complete texts after character normalization between the 50-sentence blinded sample and either the original silver training/development sets (2,150/169) or the human reference set. Of these 50 sentences, each expanded pool contains 46 in training, two in validation, and one in testing; one sentence is absent. The independent annotations support agreement analysis, but these sentences cannot provide an unseen test set for models trained on the expanded pools.

\subsection{Resource Access and Privacy Screening}
\label{sec:resource-access}
\begin{table}[!htbp]
\centering\small
\caption{Available components and supported reproduction tasks.}
\label{tab:component-access}
\begin{tabularx}{\linewidth}{@{}P{.23\linewidth}P{.30\linewidth}L@{}}
\toprule Component & Access route / available objects & Supported use and access conditions \\
\midrule
Dataset archive v0.1.3 & \href{https://doi.org/10.5281/zenodo.22698504}{Texts and labels}: source corpus, human reference (150 sentences), and Qwen/encoder Silver train/dev (2,150/169 records) & Download the named version to inspect the released labels and use the specified splits. \\
Qwen reproduction archive & \href{https://doi.org/10.5281/zenodo.22851581}{Model archive}: three Qwen LoRA adapters trained on the 2,150/169 Silver sets, selected expanded adapters, frozen predictions, prompts and scorer & Follow its model index for inference or rescoring; obtain base weights and supervision files from their respective sources. \\
Expanded pools & 9,540 adopted / 9,646 supplementary records; complete files verified locally against split manifests & Published IDs and selected outputs support inspection and rescoring. Complete inputs are available through the corresponding author for confidential peer review. \\
Additional review and historical coding & Pseudonymized texts and human layers linked in Data Availability; source hashes and cohort membership & Recalculate agreement for the documented historical annotation and assisted-review procedures. \\
Three-annotator blinded study & \href{https://github.com/AlfredJamesLi/chinese-skillspan-benchmark/tree/main/reproduction/agreement/finalguide_abc_20261003}{Agreement release}: 50 sentences, three individual label layers, source IDs, scorer and intervals & Recalculate agreement among individual annotations on the 50 formal sentences; the five familiarization and 15 practice sentences are documented separately. \\
Chinese encoder baselines & Six public \href{https://github.com/AlfredJamesLi/chinese-skillspan-benchmark/blob/0c73207c0f9b00fa87f385b78bcb7342bd00bfcb/results_snapshots/table11_evidence_20260924/HF_RELEASE.md}{Hugging Face checkpoints}; \href{https://doi.org/10.5281/zenodo.22942441}{audit DOI 10.5281/zenodo.22942441} & Frozen predictions, code and configurations support rescoring. Local rescoring verified all six runs; server checkpoint re-inference is separately recorded. \\

\bottomrule
\end{tabularx}
\par\smallskip\footnotesize\RaggedRight
Data Availability states the component-specific reuse terms. Reproduction scope follows the released materials and access conditions listed above.
\end{table}


\begin{samepage}
An automated screen examined text fields in 15 public corpus and annotation files, covering 28,635 distinct complete texts. Email, mobile-number, and landline patterns flagged 38, 74, and 50 texts, respectively; categories overlap. No such candidates were found in the human reference set or the additional review sample, but candidates occur in the Silver training/development sets and in the separate blinded agreement sample. These pattern matches may include organizational contacts and false positives and require human review. The screen does not establish the absence of other personal information or permission to redistribute texts. Aggregate counts are public; masked record-level findings remain local pending human review for privacy before release.\par
\end{samepage}

\paragraph{Limits to training reproduction.}
The JobBERT-zh fine-tuning experiments used the released Chinese recruitment encoder and its existing CRF layer. The exact original tokenizer version, some pretraining and model-selection settings, and the numbers of inputs truncated or divided into windows were not recovered. The model's input limit is 256 tokens, rather than characters; all reference spans remain in the evaluation. The exact Qwen model revision on Hugging Face was not recorded, although file checksums identify the archived files. These missing details limit reproduction of training, but the released predictions can still be used to recalculate scores.

\FloatBarrier
